\documentclass[a4paper,12pt]{article}

\usepackage[T2A]{fontenc}
\usepackage[utf8]{inputenc}
\usepackage[russian]{babel}

\usepackage{amsmath, amssymb, amsthm, mathtools}
\usepackage{helvet}
\renewcommand{\familydefault}{\sfdefault}
\usepackage{icomma}
\usepackage[left=18mm,right=18mm,top=18mm,bottom=20mm]{geometry}
\usepackage{microtype}
\usepackage{tikz}
\usepackage{tkz-euclide}
\usepackage{pgfplots}
\pgfplotsset{compat=1.18}
\usepackage{tabularx, booktabs, longtable}
\usepackage{enumitem}
\usepackage{hyperref}
\usepackage{parskip}
\usepackage{fancyhdr}
\usepackage{setspace}
\usepackage{xcolor}

\definecolor{accent}{HTML}{008080}
\definecolor{darkaccent}{HTML}{004D4D}
\definecolor{textgray}{HTML}{333333}
\definecolor{softgray}{HTML}{F2F6F6}
\definecolor{muted}{HTML}{6B7C7C}

\hypersetup{hidelinks}

\onehalfspacing
\setlength{\parindent}{0pt}
\setlength{\headheight}{14.5pt}

\pagestyle{fancy}
\fancyhf{}
\lhead{\textcolor{muted}{Логарифмические и показательные неравенства}}
\rhead{\textcolor{muted}{ЕГЭ}}
\cfoot{\thepage}

\setlist[itemize]{leftmargin=1.2em,itemsep=0.25em,topsep=0.25em}
\setlist[enumerate]{leftmargin=1.4em,itemsep=0.35em,topsep=0.35em}

\newcommand{\SectionTitle}[1]{%
  \vspace{0.8em}
  {\Large\bfseries\textcolor{darkaccent}{#1}}\par
  \vspace{0.2em}
  {\color{accent}\rule{\linewidth}{0.8pt}}
  \vspace{0.4em}
}

\newcommand{\MiniTitle}[1]{%
  \vspace{0.45em}
  {\bfseries\textcolor{darkaccent}{#1}}\par
}

\newcommand{\Important}[1]{%
  \textcolor{darkaccent}{\bfseries #1}
}

\begin{document}

\begin{titlepage}
\centering
\vspace*{28mm}

{\Large\textcolor{darkaccent}{ЕГЭ по математике}\par}
\vspace{16mm}

{\fontsize{25}{30}\selectfont\bfseries Чек-лист\par}
\vspace{5mm}
{\Large\bfseries Логарифмические и показательные неравенства\par}
\vspace{3mm}
{\large с переменным основанием, ОДЗ и методом интервалов\par}

\vspace{22mm}

{\large Лёвин Артём Александрович\par}
\vspace{2mm}
{\large эксклюзивно для Николь Саркисьянц\par}

\vspace{18mm}

{\large \today\par}

\vfill

\begin{tikzpicture}[remember picture,overlay]
  \draw[
    accent,
    line width=1.15pt
  ]
  ([xshift=22mm,yshift=28mm]current page.south west)
  .. controls
  ([xshift=62mm,yshift=52mm]current page.south west)
  and
  ([xshift=132mm,yshift=8mm]current page.south west)
  ..
  ([xshift=188mm,yshift=30mm]current page.south west);

  \draw[
    darkaccent,
    line width=0.65pt
  ]
  ([xshift=28mm,yshift=18mm]current page.south west)
  .. controls
  ([xshift=82mm,yshift=34mm]current page.south west)
  and
  ([xshift=128mm,yshift=2mm]current page.south west)
  ..
  ([xshift=178mm,yshift=20mm]current page.south west);
\end{tikzpicture}

\end{titlepage}

\SectionTitle{1. Главная идея занятия}

В логарифмических и показательных неравенствах важно не только преобразовать выражение, но и сохранить все ограничения. Если условие должно выполняться обязательно, оно записывается в системе.

\[
\text{ОДЗ и решение должны выполняться одновременно.}
\]

Поэтому в итоговом ответе берётся только та часть числовой прямой, где одновременно выполнены все условия.

\MiniTitle{Ключевой принцип}

\[
\begin{cases}
\text{ограничения},\\
\text{полученное неравенство}
\end{cases}
\quad
\Longrightarrow
\quad
\text{пересечение всех подходящих промежутков.}
\]

\SectionTitle{2. ОДЗ логарифма}

Логарифм
\[
\log_a f(x)
\]
имеет смысл только при выполнении трёх условий:
\[
\boxed{
\begin{cases}
a>0,\\
a\ne 1,\\
f(x)>0.
\end{cases}
}
\]

Если основание зависит от \(x\), то ограничения на основание тоже обязательно входят в систему.

\MiniTitle{Пример}

Для выражения
\[
\log_{2x-3}(x+2)
\]
получаем:
\[
\begin{cases}
2x-3>0,\\
2x-3\ne 1,\\
x+2>0.
\end{cases}
\]

Отсюда:
\[
\begin{cases}
x>\dfrac32,\\
x\ne 2,\\
x>-2.
\end{cases}
\]

Итоговая ОДЗ:
\[
x>\dfrac32,\qquad x\ne 2.
\]

\Important{Типичная ошибка:} писать \(2x-3\ne0\) вместо \(2x-3\ne1\). Для основания логарифма запрещено именно значение \(1\), а не только \(0\).

\SectionTitle{3. Равносильный переход для логарифмических неравенств}

Если основания одинаковые и зависят от \(x\), то удобно использовать переход:
\[
\boxed{
\log_a F(x)\ \vee\ \log_a G(x)
\Longleftrightarrow
\begin{cases}
a>0,\\
a\ne1,\\
F(x)>0,\\
G(x)>0,\\
(a-1)\bigl(F(x)-G(x)\bigr)\ \vee\ 0.
\end{cases}
}
\]
Знак \(\vee\) означает любой знак сравнения:
\[
>,\quad <,\quad \ge,\quad \le.
\]

\MiniTitle{Почему появляется множитель \(a-1\)}

Если \(a>1\), логарифмическая функция возрастает.

Если \(0<a<1\), логарифмическая функция убывает.

Множитель \(a-1\) автоматически учитывает оба случая.

\MiniTitle{Пример}

Решим неравенство:
\[
\log_x(x+1)<\log_x4.
\]

Пишем систему:
\[
\begin{cases}
x>0,\\
x\ne1,\\
x+1>0,\\
4>0,\\
(x-1)\bigl((x+1)-4\bigr)<0.
\end{cases}
\]

Последнее неравенство:
\[
(x-1)(x-3)<0.
\]

Отсюда:
\[
1<x<3.
\]

Ответ:
\[
\boxed{x\in(1;3)}.
\]

\SectionTitle{4. Система, совокупность и пересечение штриховок}

Ограничения всегда записываются под знаком системы, потому что они должны выполняться одновременно.

Если внутри решения появляется совокупность, то она отвечает за выбор одного из возможных случаев. Но вся совокупность всё равно может находиться внутри общей системы ограничений.

\[
\begin{cases}
x>-1,\\
\left[
\begin{array}{l}
x>3,\\
x<-5
\end{array}
\right.
\end{cases}
\]

Здесь нужно взять только те значения, которые удовлетворяют и ограничению \(x>-1\), и одному из условий совокупности.

Итог:
\[
x\in(3;+\infty).
\]

\MiniTitle{Смысл на числовой прямой}

\begin{center}
\begin{tikzpicture}[x=0.9cm,y=0.9cm]
  \draw[->,line width=0.8pt] (-6.5,0) -- (5.5,0);
  \foreach \x/\lab in {-5/{-5},-1/{-1},3/{3}}
  {
    \draw (\x,0.08) -- (\x,-0.08);
    \node[below] at (\x,-0.08) {\(\lab\)};
  }
  \draw[accent,line width=2pt] (-1,0.35) -- (5,0.35);
  \draw[darkaccent,line width=2pt] (3,0.72) -- (5,0.72);
  \draw[darkaccent,line width=2pt] (-6,0.72) -- (-5,0.72);
  \draw[fill=white,line width=0.8pt] (-1,0.35) circle (0.08);
  \draw[fill=white,line width=0.8pt] (3,0.72) circle (0.08);
  \draw[fill=white,line width=0.8pt] (-5,0.72) circle (0.08);
  \node[right,textgray] at (5.1,0.35) {\small \(x>-1\)};
  \node[right,textgray] at (5.1,0.72) {\small \(x>3\)};
  \node[above,textgray] at (4.1,1.0) {\small общая область: \((3;+\infty)\)};
\end{tikzpicture}
\end{center}

\SectionTitle{5. Метод интервалов для дробей и произведений}

Если нужно решить неравенство вида
\[
\frac{P(x)}{Q(x)}\ \vee\ 0,
\]
то на одной числовой прямой отмечаются:

\begin{itemize}
\item нули числителя;
\item нули знаменателя;
\item точки, запрещённые ОДЗ.
\end{itemize}

Нули знаменателя всегда выколоты.

\MiniTitle{Алгоритм}

\begin{enumerate}
\item Перенесите всё в одну сторону.
\item Приведите выражение к произведению или дроби.
\item Найдите нули всех множителей.
\item Отметьте точки на одной числовой прямой.
\item Расставьте знаки на промежутках.
\item Выберите промежутки, подходящие под знак неравенства.
\item Пересеките результат с ОДЗ.
\end{enumerate}

\MiniTitle{Важное правило}

Если множитель встречается чётное число раз, знак при переходе через его корень не меняется.

Например:
\[
(x-1)^2\ge0
\]
выполняется при всех \(x\in\mathbb{R}\), потому что квадрат неотрицателен.

\SectionTitle{6. Нельзя выбрасывать знаменатель}

Если в решении появляется дробь, нельзя писать условие только для числителя.

Ошибочная логика:
\[
\frac{P(x)}{Q(x)}>0
\quad\Longrightarrow\quad
P(x)>0.
\]

Правильно:
\[
\frac{P(x)}{Q(x)}>0
\]
решается методом интервалов с учётом и числителя, и знаменателя.

\MiniTitle{Почему так}

Знак дроби зависит от двух частей:
\[
\operatorname{sign}\frac{P(x)}{Q(x)}
=
\operatorname{sign}P(x)\cdot\operatorname{sign}Q(x).
\]

Поэтому знаменатель нельзя просто убрать.

\SectionTitle{7. Показательные неравенства с постоянным основанием}

Если основание постоянно, то правило зависит от величины основания.

\[
a^{f(x)}\ \vee\ a^{g(x)},\qquad a>0,\quad a\ne1.
\]

Если \(a>1\), знак сохраняется:
\[
a^{f(x)}\le a^{g(x)}
\Longleftrightarrow
f(x)\le g(x).
\]

Если \(0<a<1\), знак меняется:
\[
a^{f(x)}\le a^{g(x)}
\Longleftrightarrow
f(x)\ge g(x).
\]

\MiniTitle{Пример}

\[
\frac{6\cdot3^x}{2^x}+\frac{3^x}{2^x}\le10{,}5.
\]

Выносим общее основание:
\[
6\left(\frac32\right)^x+\left(\frac32\right)^x\le10{,}5.
\]

\[
7\left(\frac32\right)^x\le\frac{21}{2}.
\]

\[
\left(\frac32\right)^x\le\frac32.
\]

Так как \(\dfrac32>1\), знак сохраняется:
\[
x\le1.
\]

Ответ:
\[
\boxed{x\in(-\infty;1]}.
\]

\SectionTitle{8. Показательные неравенства с переменным основанием}

Если основание зависит от \(x\), полезен переход:
\[
\boxed{
F(x)^{u(x)}\ \vee\ F(x)^{v(x)}
\Longleftrightarrow
\begin{cases}
F(x)>0,\\
\bigl(F(x)-1\bigr)\bigl(u(x)-v(x)\bigr)\ \vee\ 0.
\end{cases}
}
\]

Если используется школьная схема сравнения показателей как у показательной функции, отдельно контролируйте ситуацию \(F(x)=1\). Для строгих неравенств такие точки обычно автоматически не подходят. Для нестрогих их лучше проверить отдельно.

\MiniTitle{Если справа стоит единица}

Единицу удобно представить как нулевую степень:
\[
1=F(x)^0.
\]

Тогда:
\[
F(x)^{u(x)}>1
\Longleftrightarrow
F(x)^{u(x)}>F(x)^0.
\]

После этого:
\[
\begin{cases}
F(x)>0,\\
\bigl(F(x)-1\bigr)u(x)>0.
\end{cases}
\]

\MiniTitle{Пример}

\[
(x^2-4x+4)^{x^2-x-6}>1.
\]

Здесь:
\[
F(x)=x^2-4x+4=(x-2)^2.
\]

Ограничение:
\[
(x-2)^2>0
\quad\Longleftrightarrow\quad
x\ne2.
\]

Переход:
\[
\bigl((x-2)^2-1\bigr)(x^2-x-6)>0.
\]

Раскладываем:
\[
(x-1)(x-3)(x-3)(x+2)>0.
\]

То есть:
\[
(x-1)(x-3)^2(x+2)>0.
\]

С учётом \(x\ne2\) получаем:
\[
\boxed{x\in(-\infty;-2)\cup(1;2)\cup(2;3)\cup(3;+\infty)}.
\]

\SectionTitle{9. Модуль в основании}

Если основание содержит модуль, сначала выписываются ограничения.

Для выражения
\[
\log_{|x-1|}(x+2)
\]
получаем:
\[
\begin{cases}
|x-1|>0,\\
|x-1|\ne1,\\
x+2>0.
\end{cases}
\]

Разбираем ограничения:
\[
|x-1|>0
\Longleftrightarrow
x\ne1.
\]

\[
|x-1|\ne1
\Longleftrightarrow
\begin{cases}
x-1\ne1,\\
x-1\ne-1.
\end{cases}
\]

То есть:
\[
x\ne2,\qquad x\ne0.
\]

Итог:
\[
x>-2,\qquad x\ne0,\qquad x\ne1,\qquad x\ne2.
\]

\SectionTitle{10. Корень в степени}

Если в показателе есть корень чётной степени, обязательно учитывается подкоренное выражение.

\[
\sqrt{x+2}
\]
имеет смысл при:
\[
x+2\ge0.
\]

То есть:
\[
x\ge-2.
\]

\MiniTitle{Равенство с корнем}

Для уравнения
\[
\sqrt{F(x)}=G(x)
\]
безопасный равносильный переход:
\[
\boxed{
\begin{cases}
G(x)\ge0,\\
F(x)=G^2(x).
\end{cases}
}
\]

Это важно, потому что после возведения в квадрат могут появиться лишние корни.

\MiniTitle{Пример}

\[
2\sqrt{x+2}=x+1.
\]

Пишем:
\[
\begin{cases}
x+1\ge0,\\
4(x+2)=(x+1)^2.
\end{cases}
\]

Второе уравнение:
\[
4x+8=x^2+2x+1.
\]

\[
x^2-2x-7=0.
\]

\[
x=1\pm2\sqrt2.
\]

С учётом \(x+1\ge0\) подходит только:
\[
x=1+2\sqrt2.
\]

\SectionTitle{11. Типичные ошибки}

\begin{enumerate}
\item Забыть записать ОДЗ логарифма.
\item Написать для основания логарифма \(a\ne0\), но забыть \(a\ne1\).
\item Решить ограничения отдельно и не пересечь их с ответом.
\item Заменить систему совокупностью там, где условия должны выполняться одновременно.
\item В дробном неравенстве смотреть только на числитель.
\item Отметить нуль знаменателя закрашенной точкой.
\item Не заметить чётную кратность корня.
\item При переменном основании не учесть множитель \(F(x)-1\).
\item В неравенстве с корнем забыть условие подкоренного выражения.
\item После возведения в квадрат не проверить условие на правую часть.
\end{enumerate}

\SectionTitle{12. Короткий алгоритм решения}

\begin{enumerate}
\item Определите тип неравенства: логарифмическое, показательное, дробное, с модулем, с корнем.
\item Выпишите все ограничения.
\item Примените подходящий равносильный переход.
\item Приведите выражение к произведению или дроби.
\item Решите методом интервалов.
\item Пересеките найденное решение со всеми ограничениями.
\item Проверьте особые точки: выколотые точки, корни знаменателя, точки \(F(x)=1\), границы ОДЗ.
\end{enumerate}

\newpage

\SectionTitle{Тренировочный блок}

Решите неравенства. Ответ запишите в виде промежутков.

\begin{enumerate}[label=\textbf{\arabic*.}]
\item
\[
\log_x(x+1)<\log_x4.
\]

\item
\[
\log_{x-2}(x+1)\ge\log_{x-2}5.
\]

\item
\[
\frac{6\cdot3^x}{2^x}+\frac{3^x}{2^x}\le10{,}5.
\]

\item
\[
\log_{x+1}(x+3)>\log_{x+1}(2x).
\]

\item
\[
\log_{2x-3}(x+2)\le\log_{2x-3}(7-x).
\]

\item
\[
\begin{cases}
x>0,\\
x\ne1,\\
x>\dfrac75,\\
\dfrac{(2x-1)(x-3)}{x-1}\ge0.
\end{cases}
\]

\item
\[
\log_x\dfrac{x+2}{x-1}\le\log_x3.
\]

\item
\[
(x^2-4x+4)^{x^2-x-6}>1.
\]

\item
\[
\log_{|x-1|}(x+2)>\log_{|x-1|}4.
\]

\item
\[
|x-1|^{2\sqrt{x+2}}<|x-1|^{x+1}.
\]
\end{enumerate}

\newpage

\SectionTitle{Ответы для самопроверки}

\textcolor{muted}{Ключ-пароль для открытия блока: \(\text{NIKOL-LOG-EGE}\).}

\vspace{0.6em}

\renewcommand{\arraystretch}{1.45}
\begin{table}[h]
\centering
\caption{Ответы к тренировочному блоку}
\begin{tabularx}{\linewidth}{@{}cX@{}}
\toprule
\textbf{№} & \textbf{Ответ} \\
\midrule
1 & \((1;3)\) \\
2 & \((2;3)\cup[4;+\infty)\) \\
3 & \((-\infty;1]\) \\
4 & \((0;3)\) \\
5 & \(\left(2;\dfrac52\right]\) \\
6 & \([3;+\infty)\) \\
7 & \(\left[\dfrac52;+\infty\right)\) \\
8 & \((-\infty;-2)\cup(1;2)\cup(2;3)\cup(3;+\infty)\) \\
9 & \((0;1)\cup(1;2)\cup(2;+\infty)\) \\
10 & \((0;1)\cup(1;2)\cup(1+2\sqrt2;+\infty)\) \\
\bottomrule
\end{tabularx}
\end{table}

\end{document}